% Dataset descriptions and licenses (App. \ref{app:data}).  Facts from the official
% release pages and dataset papers (checked 2026-09-30); our processing from
% tools/unified/readers/*.py.
\paragraph{Datasets.} All data are public releases, recorded by their authors with ethics
approval and informed consent; we collected no new data. Table~\ref{tab:licenses} lists every
release, the version we used and its license. ``Our use'' states what we changed from the
release.

\emph{KU Leuven} \citep{das2016kul,biesmans2017auditory}. 16 normal-hearing listeners,
64-channel BioSemi ActiveTwo, released at 128\,Hz. Four Dutch short stories read by male
narrators, played at 60\,dB(A) (A-weighted sound level) through insert phones, either dichotically
(dry) or filtered with head-related transfer functions (HRTFs) to $\pm90^\circ$. 20 trials per listener (8 of about 6\,min and 12 two-minute
repetitions, about 72\,min); the attended ear is fixed within a trial. Our use: all 20
trials, envelopes from the dry story recordings, content groups = the four stories.

\emph{DTU} \citep{fuglsang2017dtu}. 18 normal-hearing listeners, 64-channel BioSemi at
512\,Hz. Danish fictional stories read by a male and a female professional speaker, about
50\,s per trial, 60 trials per listener, talkers at $\pm60^\circ$ through simulated room
impulse responses (anechoic, mildly and highly reverberant rooms) and insert earphones; the
target talker is cued before each trial. Our use: the two-talker trials only (single-talker
trials and one trial that stops after 19.6\,s are dropped).

\emph{Fuglsang'20} \citep{fuglsang2020effects}. 44 listeners aged 51--76, 22 with
sensorineural hearing loss and 22 with normal hearing; 64-channel BioSemi at 512\,Hz. Two
audiobooks read by a male and a female talker at $\pm90^\circ$ (non-individualised HRTFs,
insert phones), 48 trials of about 50\,s, of which 32 are two-talker; the competing talker
starts 3--5\,s after the target. The release contains no audio, only the authors'
envelopes. Our use: the two-talker trials (32 per listener; 25 for one listener) and the provided envelopes of the
unamplified stimuli (hearing-impaired listeners heard an amplified version), so every listener
shares one envelope per stimulus.

\emph{ESAA} \citep{li2022esaa}. The release holds 17 normal-hearing listeners (the paper
reports 20), 64-channel EEG recorded at 1\,kHz (BrainProducts files). Mandarin short
stories told by a female and a male professional storyteller, HRTF-filtered to $\pm90^\circ$,
32 trials of 55--90\,s per listener (one recording ends after 16\,s). Our use: all trials, including those the authors flag
for poor comprehension answers; the stimulus reaches the ear about 0.3\,s after the start
marker, which we correct (App.~\ref{app:data}, alignment checks).

\emph{KUL-255ch} \citep{narayanan2021uhd}. 30 normal-hearing male listeners aged 22--35,
255-electrode 10-5 cap (253 EEG channels in the files), Compumedics SynAmps RT at 1\,kHz.
Two parts of two of the KU Leuven Dutch stories, HRTF-filtered to the left and right, four
trials of about 6.5\,min, one per attended ear and story part. Our use: all four trials.
Because the stories are those of KU Leuven, the two datasets are never mixed across training
and test. The authors corrected their stimulus table in August 2026; our attended labels are
checked on the data, as the labelled attended envelope is the one the EEG tracks
(Table~\ref{tab:main}).

\emph{Broderick} \citep{broderick2018semantic}. 33 native English listeners (34 recorded),
128-channel BioSemi ActiveTwo plus mastoids at 512\,Hz, released at 128\,Hz. Two Jules Verne
novels read by two male speakers, one per ear over headphones (\emph{Twenty Thousand Leagues
Under the Sea} left, \emph{Journey to the Centre of the Earth} right), 30 trials of about
60\,s; each listener attends the same side in every trial (listeners 1--17 left, 18--33
right). The release contains envelopes, not audio. Our use: the ``Cocktail Party'' part,
the provided envelopes, and the 64 channels nearest the canonical layout (all 128 in
Table~\ref{tab:channels}).

\emph{SparrKULee} \citep{accou2024sparrkulee}. 85 normal-hearing native Flemish listeners,
64-channel BioSemi ActiveTwo (recorded at 8192\,Hz, distributed at 1024\,Hz). Single-talker
Flemish audiobooks and podcasts, 6--10 trials of about 15\,min per listener. Our use: the
recordings of the training set of the auditory-EEG challenge at the 2024 IEEE International
Conference on Acoustics, Speech and Signal Processing (ICASSP); we drop the stories that
also occur in KU Leuven and the two speech-in-noise runs, and re-time the audio to the EEG
clock with the recorded trigger pulses.

\emph{MEG-MASC} \citep{gwilliams2023megmasc}. 27 English speakers, 208-channel KIT axial
gradiometer MEG at 1\,kHz. Four synthesised English fictional stories (about 2\,h, one or two
sessions), interleaved with short word lists, with comprehension questions. Our use: every
played audio file as one trial, sound onsets corrected by the delay measured on the recorded
audio channel, and all 208 sensors, read by a 208-channel input layer. A per-listener
principal component analysis (PCA) to 64 components, which we used first, fits different components for each listener; it
matched the 208-sensor model on the content folds but cannot transfer to an unseen listener
(\losoMEGpca\% instead of \losoMEGFive\% under leave-one-listener-out evaluation), so no model in the paper uses it.

\emph{MAESTRO} \citep{maestro2026}. 16 listeners aged 18--30, 32-channel ANT Neuro eego at
500\,Hz, plus eye gaze, pupil, egocentric video and head motion (not used here). English
speech from six loudspeakers in a room: four attendable talkers and two distractors, with
background noise; 100 evaluation trials of about 30\,s and 5 familiarisation trials. Our use:
the EEG of the 100 evaluation trials, with the synchronisation described below.

\emph{Excluded.} NJU \citep{zhang2023learnable} and AV-GC \citep{rotaru2024what}, for the
reasons given below; their licenses are in Table~\ref{tab:licenses}.

\paragraph{Licenses.} Every license permits our use, which is non-commercial research. We
redistribute no data, stimuli or derived features; the released code contains no data and
reads the original files, so the share-alike terms (KU Leuven, MAESTRO) do not apply to it. SparrKULee's stimuli
are licensed for non-commercial use only, and 107 of its recordings are available only on
request; we used the recordings distributed for the ICASSP 2024 challenge.
\begin{table*}[h]\centering\small
\caption{Releases, versions and licenses (checked on the release pages, September 2026).
CC: Creative Commons; BY: attribution; NC: non-commercial; SA: share-alike; CC0: public-domain
dedication. BIDS: Brain Imaging Data Structure; RDR: Research Data Repository; OSF: Open
Science Framework. Excluded datasets are listed for completeness.}
\label{tab:licenses}
\resizebox{\textwidth}{!}{%
\begin{tabular}{llllll}\toprule
Dataset & repository & identifier & version used & license & access\\\midrule
KU Leuven & Zenodo & 10.5281/zenodo.4004271 & v2.0 (2019)$^a$ & CC BY-NC-SA 4.0 & open\\
DTU & Zenodo & 10.5281/zenodo.1199011 & v1 (2018) & CC BY-NC 4.0 & open\\
Fuglsang'20 & Zenodo & 10.5281/zenodo.3618205 & 2020 (BIDS 1.3.0) & CC BY 4.0 & open (audio on request)\\
ESAA & Zenodo & 10.5281/zenodo.7795585 & v1.0 (2022) & CC BY 4.0 & open\\
KUL-255ch & Zenodo & 10.5281/zenodo.4518754 & v1.1.0 (2021) & CC BY-NC 4.0 & open\\
Broderick & Dryad & 10.5061/dryad.070jc & v3 (2019) & CC0 1.0 & open\\
SparrKULee & KU Leuven RDR & 10.48804/K3VSND & v3.1 (2023) & CC BY-NC 4.0 & open; 107 recordings on request\\
MEG-MASC & OSF & 10.17605/OSF.IO/AG3KJ & 2022 & CC0 1.0 & open\\
MAESTRO & Hugging Face & aspire-osu/maestro-eeg-dataset & 2026 & CC BY-NC-SA 4.0 & open\\\midrule
NJU (excl.) & Zenodo & 10.5281/zenodo.7253438 & v1 (2022) & CC BY 4.0 & open\\
AV-GC (excl.) & Zenodo & 10.5281/zenodo.11058711 & v1 (2024) & CC BY-SA 4.0 & open (13 of 16 listeners)\\
\bottomrule\end{tabular}}
\par\smallskip\raggedright\footnotesize $^a$Version 2.1 (July 2026, 10.5281/zenodo.21415059)
corrects a 70-ms misalignment between EEG and stimuli, after our experiments. Our lag window
spans about 1\,s and its start was chosen on validation data (125\,ms, consistent with the
offset scan of our bundle, which peaks one 125-ms step after 0\,s); we have not re-run the
experiments on version 2.1.
\end{table*}
